% latexmk -cd -lualatex works/L_2/appendix.tex
% latexmk -c -cd works/L_2/appendix.tex

\documentclass{ees}

\pghyphenation{english}{}

\newlist{lyricslist}{description}{1}
\setlist[lyricslist]{
  partopsep=0pt,
  labelindent=0em,
  labelwidth=7em,
  leftmargin=7em,
  labelsep=0pt,
  first=\ltseries,
  font=\normalfont\itshape\ltseries,
  style=multiline
}

\newenvironment{lyrics}[1]{%
  \subsection{#1}\nopagebreak%
  \begin{lyricslist}%
  \let\voice\item%
}{%
  \end{lyricslist}%
}

\begin{document}

\clearscrheadfoot

\chapter{Prefatory pages}

\textit{[bc part, fol. 1r]}\\
Zwey neue | und extra luſtige muſicaliſche | Tafel-Stücke. | I. | Der Wienneriſche Tändlmarckt. | Mit 4. Singſtimmen, 2. Violinen und Baſſo ordinario. | Die Singſtim- | men ſtellen vor: Den Sägfeiller, Hollhippen-Kramer, einen Savoyarden | mit ſeinem Spielwerck, und einen Marcktſchreyer. Die Inſtru- | menten aber imitiren das ſogenannte Winne- | riſche Breinglöckl. | II. | Die Bauren-Richters-Wahl. | Mit 5. Singſtimmen. 2. Violinen und Baſſo ordinario. Die Singer ſind: | Der Bartlme Zimmermann, als Senior, der Paul Schnepfendreck, der | Böhmiſche Mazko, der Schwäbiſche Wirth, und der einäugige | Schulmeiſter, welche ein jeder in ſeiner Sprache, | die Vota zur Wahl geben. | Beyde Stücke/ ſind nach Art des Quodlibete eingerich- | tet, und jedes, zur angenehmen Abwechslung, mit 7. Arien, und darzu | gehörigen Recitativen verſeheh, mithin bey vornehmen Gaſt- | mahlen und anderen luſtigen Geſellſchaften mit | groſſem Applauſu zu produciren. | Von dem Autore des Muſicaliſchen Calenders | Gregorio Joſepho Werner, | Hochfürſtlich Eſtorhaſiſchen Capellmeiſtern zu Eiſenstadt in Ungarn. | AUGSPURG, | Verlegt von Johann Jacob Lotters ſeel. Erben.

\textit{[fol. 2r]}\\
Avertiſſement. | Hier werden einer courieuſen Welt/ zur Luſt und | Zeit-vertreiblichen Ergötzlichkeit zwey Muſicali- | ſche Quodlibetica vorgeſtellet. Daß erſtere | betitult ſich/ der Wieneriſche Tändelmarckht/ | allwo in dem Choro alles natürlich begriffen/ | wie jedwe- | dere ihre Waaren allda zum Verkauf pflegen auszuruffen/ | worunter das ſogenannte Breinglöckel in denen Inſtru- | mental-Stimmen Imitations-Weiß ſich hören läßt. Als- | dann erfolget Recitativ- und Arien-mäßig der Sägfeil- | ler/ Savojard, Hollhippen-Kramer/ ſamt dem Arzt | auf der Freyung. Alles iſt leicht zu tractiren/ | anbey mit

\textit{[fol. 2v]}\\
guten Fleiß kurz verfaſſet, damt das Auditorium immer | eine angenehme Veränderung hören möge. Das zweyte | Stuck/ wird nach allhießiger Landsart benamßet die Bon- | thätung/ oder die Bauren-Richters-Wahl/ ſo beſtehet | aus 5. Perſonagen; als der Barthlme Zimmermann/ | der Paul Schnepfendreck/ der Schwabnwürth/ der | Böhmiſch Mazko, und der Schulmeiſter/ welcher mit | ſeiner Schul-Jugend zu buchſtabieren anfanget. Alsdann | wird zur Wahl geſchritten/ mit einer vorgehenden Bauren- | Oration, da endlich nach langem Wortſtreit und Knittel- | haften Expreſſionen der Richter einſtimmig erwählet wird. | In allen beyden Pieçen wird ſodann jedem Vocaliſten/ | in ſeiner Perſon hauptſächlich die Pronunciation wohl | anrecommendiret. | Der ich mich übrigens zu der re- | ſpectivè geneigten Liebhaberen Dienſte/ | fernerhin gehor- | ſamſt anerbietend empfehle. | Ein befliſſener Diener | Gregorius Joſeph Werner.

\cleardoublepage
\chapter{Lyrics}

\textit{Roles:} Der Schullmeiſter (T) · Der Schwabnwürth (T) · Paul Schnepfendrekh (T) · \\Der Bömiſch Mazko (B) · Barthlme Zimmermann (B)
\enlargethispage{2\baselineskip}

\subsection{1 Introitus}

\begin{lyrics}{2 Recitativo}
  \voice[Schullmeiſter]
  Kinder lährnts bräf drauf\\
  und ſagts mär bald auf,\\
  I kan enkh nit länger mehr daſtehn\\
  I muß auf dö Bohnthätung gehn.
\end{lyrics}

\begin{lyrics}{3 Chorus}
  \voice[alle]
  Cas ti ri bom.
\end{lyrics}

\begin{lyrics}{4 Recitativo}
  \voice[Schullmeiſter]
  Allo ſei ſtät und gmach,\\
  ſingts mir hiezt alli nach.
\end{lyrics}

\begin{lyrics}{5 Chorus}
  \voice[alle]
  Änmal äns iſt äns\\
  ä Nulla daß iſt käns\\
  zwämahl 2 iſt vieri\\
  zwämahl 4 iſt ächti\\
  zwämahl 5 iſt zöhni\\
  zwämahl 10 iſt zwänzg\\
  zöhnmahl 10 iſt hundert\\
  zöhnmahl hundert iſt tauſend.
\end{lyrics}

\begin{lyrics}{6 Recitativo}
  \voice[Barthlme]
  Obwohl obſchon\\
  obgleich und ſintemall\\
  der leydig krach düerre Tod\\
  mit ſeinem Pfeil\\
  in ſcheller Eyl\\
  unſern Richter den Thoma Pökh\\
  den alten Gekh\\
  in oänen hui hat troffen,\\
  und Er äß wie ä Scheer\\
  in dErden eini gſchloffen,\\
  ſo wär uns do zum Troſt\\
  ſei höniſüeſſer Moſt\\
  zu oänen Labſall\\
  auf ſeinem Todtenmahl.\\
  Nu brauch mär halt än Man\\
  der uns recht mit Verſtand\\
  khan gehen an dö Hand\\
  und unſrä gantzen Gmoän\\
  ganz muädäſeel alloän\\
  recht gſcheidt vorſtehen khan.

  \voice[Schullmeiſter]
  Warhaffti warhaffti\\
  dö Röd iß ſaffti,\\
  der Parthlme Zimmermann\\
  halt gar äwei ſi Sermonion,\\
  der Text iſt tröffli nobl\\
  är ſchmäkt mei Seeil nit nahn Hobl.
\end{lyrics}

\begin{lyrics}{7 Aria 1}
  \voice[Barthlme]
  Ein Zimmermann der etwas khan\\
  iſt allerdings zu loben,\\
  weil er ſich plagt ſein Leben wagt\\
  biß an die Spize oben,\\
  was denckt er nicht dicht her und hin\\
  biß er daß Holtz außzirckelt\\
  da geht die edle Zeit dahin,\\
  bevor wann er waß ſchnirkelt,\\
  än ſolcher Mann ſtund uns wohl an\\
  wann er ſolt Richter weren,\\
  er iſt än Würth dem Dorf än Zürd,\\
  den könt mär billich ehren.
\end{lyrics}

\begin{lyrics}{8 Recitativo}
  \voice[Schnepfendrekh]
  Nä nä auß deiner Wahl wird nit viel drauß,\\
  der Barthlme hats hintern Ohrn\\
  er ſticht uns d’Narren hint und vorn,\\
  äs wurd bald hoäſſen bloß märn Hobl auß.

  \voice[Schwabnwürth]
  Der Bruäder Schnepfädreckh hat recht,\\
  Er iſcht von oänen gſcheiden Gſchlecht\\
  komt offt zum Tiſch der groſſä Herrä,\\
  dort khän er kluge Sächä hörä.

  \voice[Mazko]
  Mär hamär ä ſo Sinn\\
  und kemär nit mit Wahl auf ihm\\
  mir könnä ja Schwabnwürth nähmä\\
  vor Richter wohl erkhennä,\\
  er liebt ja ä dö Gläßl Wein\\
  da kamär Gmoän bräf luſti ſeyn.
\end{lyrics}

\begin{lyrics}{9 Aria 2}
  \voice[Mazko]
  Bey Schwabnwürth bey Schwabnwürth\\
  da hed mär Dorf ä Zürd,\\
  wär ſchon Mann der was khan\\
  die Sachä wohl ausführt,\\
  hamär ghert zu Inglſtadt\\
  wie auf den Gretſcheri hat\\
  ſtadli wohl ſtudirt.
\end{lyrics}

\begin{lyrics}{10 Recitativo}
  \voice[Barthlme]
  Bakts önk mit dem Schwäwerlegaun\\
  könts ja umb oänen andern ſchaun,\\
  wiſts denn nit daß dSuppen Schwabn\\
  müſſen s’lözti Wort all habn.

  \voice[Schullmeiſter]
  Herr Barthlme hat recht,\\
  der wär mir ä zu ſchlecht.

  \voice[Schnepfendrekh]
  Mär ſagt ja daß er tröffli wohl\\
  daß Hüener Buckhä khönä ſoll,\\
  mir löſſten vor die Oär bräf Geld,\\
  ſo wär dö Gmoän ja gut beſtellt.

  \voice[Barthlme]
  Äs ſcheint wohl daſt än ſtinkäde Nahmä haſt,\\
  geh lad die ei bey öem zu Gaſt,\\
  du kanſt dö Oär all ſelbſt verzöhrn,\\
  äs wird dirs ä von unß gwiß koäner wöhrn.
\end{lyrics}

\begin{lyrics}{11 Aria 3}
  \voice[Barthlme]
  I mit mein Verſtands Maß-Stab\\
  möſſe alle Mängl ab\\
  mi ſoll koäner nit anführn\\
  hab no Roſämiſch in Hürn\\
  hab Roſomiſch in Hirn\\
  wölts nit folgen laſt äs ſeyn,\\
  I gib koäni Wahl mehr drein.
\end{lyrics}

\clearpage
\begin{lyrics}{12 Recitativo}
  \voice[Schullmeiſter]
  Mei Moänung bleibt halt wie zuvor\\
  daß man den Parthl alß Senior\\
  zum Richter Ambt erwöhlet\\
  ſo wär als guet beſtellet.

  \voice[Schwabnwürth]
  Mir brauchä ja än höfflichä Männ\\
  der jeden freundlich thuä khän\\
  der Parthl iſcht von oän widrig geſinten humor,\\
  er ſchlueg uns die Bandhäkä zwiſchä daß Ohr.

  \voice[Schnepfendrekh]
  Mei Wahl dö bleibt ſchon bey dem Schwabn\\
  da könnä mär gute Schnöäka habn,

  \voice[Barthlme]
  zumall wann ers aus ſein Oärn ausbrüeth\\
  und s’hundert vor än Bazen bieth.

  \voice[Schullmeiſter]
  Mit dem iß gar ä gfäl die Sach,\\
  dann läffet ihm ä Haaß nur nach\\
  ſo thät er ſi glei verkriechen,\\
  mir könten von dem gleichwoll was anders rüechen.

  \voice[Schnepfendrekh]
  Öß ſeits mär alli wie es ſcheint\\
  dem armen Schwäwerle mechti feund.
\end{lyrics}

\begin{lyrics}{13 Aria 4}
  \voice[Schnepfendrekh]
  Ein Schwab mueß ſchwören wie bekandt\\
  wann er will wandern auß dem Land,\\
  das lözte Wort mueß bey ihm ſeyn,\\
  dö Sprach ſoll er auch halten rein,\\
  zurukh ſoll er auch wider khern,\\
  ſo thoäns dö Landsleuth all begehrn.
\end{lyrics}

\begin{lyrics}{14 Recitativo}
  \voice[Mazko]
  Hamär ſchon bald finſter Nacht,\\
  iſt mär no koä Richter gmacht.

  \voice[Schwabnwürth]
  So wärthet halt biß morgä früe\\
  biß der Halter treibt die Küeh.

  \voice[Barthlme]
  I kim enk nit bin wie ä Bock.

  \voice[alle andern]
  So legt di halt dö gantze Gmoän in Stokh.
\end{lyrics}

\subsection{15 Nacht-Stuck}

\begin{lyrics}{16 Aria 5}
  \voice[Schwabnwürth]
  Auf auf ihr Völkle es täget ſchon,\\
  I hör die Engle ſinge\\
  daß Hänle kreht die Sonn aufgeht,\\
  die Geiſchbökh frölich ſpringe,\\
  ſo wölln wir gaun zur Wähöe ſchaun,\\
  damit ſie wohl und guet gelinge.
\end{lyrics}

\begin{lyrics}{17 Recitativo}
  \voice[Barthlme]
  Öß Dalkhen, Limeln und öß Narren!\\
  wanß lieber thäts dö Zeit erſpahren\\
  und bald oän Richter wöhln.

  \voice[Schnepfendrekh]
  Döß iß mär wohl ä Zimmerſtokh\\
  von zöhentauſend Ein.

  \voice[Schullmeiſter]
  Er ſteht ämall alß Senior\\
  der ganzen Gmoän anjezo vor.

  \voice[Mazko]
  Parthl iſt mär wie grob Holz,\\
  daß mues mit Hackhen klieben\\
  moän ſchier är hat än Richtärſtolz\\
  will ſie mit Gwald eiſchieben.

  \voice[Barthlme]
  Mazko Mazko\\
  ti dabal fazgo

  \voice[Mazko]
  Du will mir gar Ohrfeign gebn\\
  ſchwör I dir bey Leib und Löbn\\
  muſt über Schwänzär ſpringä\\
  und s’grimmi Tod bald ſinga.

  \voice[Barthlme]
  Du wirſt ja wohl koä Ziska ſeyn.

  \voice[Mazko]
  Glei halt Maul, bin I nit faul\\
  ſchlag mär wie ein Trumi drein,\\
  wann mahl bemiſch Kopp auffez\\
  I ſchon alß zu Boden fez.\\
  Schneeberg mueß ſi vor mir buckhn,\\
  Ötſcher hau I gar in Stuckhn;\\
  Monſchein züttert und ſi verkriecht\\
  wann er Mazko zorni ſieht.

  \voice[Schwabnwürth]
  Seyd doch friedlich mächt ein Ende,\\
  ich geb mein Wähl behende\\
  auf den Paula Schnepfendrekh.

  \voice[Mazko]
  Nehmärn auf mein Wort nur köck.

  \voice[Schullmeiſter]
  Wanns den anderſt nit khan ſeyn\\
  gib I mein Willen drein.

  \voice[Barthlme]
  Öß ſeyts Leuth äß is ä Grauß\\
  bloſts mär alle den Hobl auß.
\end{lyrics}

\begin{lyrics}{18 Aria 6}
  \voice[Barthlme]
  Äß iß hiez beym Bauren faſt gar koä Reſpect\\
  wann oäner viel Jahr ſchon bey öenen erſtreckt,\\
  will lieber mi hiezund in d’Stadt eini ſezen,\\
  dorth woäß mär dö Zimmerleuth böſſer zu ſchezen.
\end{lyrics}

\begin{lyrics}{19 Recitativo}
  \voice[Schnepfendrekh]
  Wälß mi den all zum Richter wolt,\\
  ſo merkh I daß I folgen ſolt,\\
  der Schullfur ſoll an meiner ſtat\\
  ſchön dancken allen vor dö Gnad.

  \voice[Schullmeiſter]
  Sonderbahre \\
  wohl achtbare!\\
  auch Knödlvöſte.\\
  I dankh enk aufs allerböſte\\
  vor gehabte Müeh und Sorgen,\\
  daß ihr uns habt dieſen Morgen\\
  ſtats verſtorbnen Thoma Bökh\\
  unſern Paulen Schnepfendrekh\\
  ganz oänhelligli erwöhlet\\
  wohl beſtöllet,\\
  daß er ſoll döß gantze Orth\\
  fort und fort\\
  weißli regiern,\\
  und unß allen dö Gurgl ſchmirn.
\end{lyrics}

\begin{lyrics}{20 Aria}
  \voice[Schullmeiſter]
  Mi braucht man zu allen Sachen\\
  waß in Dorf nur immer iſt,\\
  mues als gleich und öwen machen,\\
  wans ſcho wär ä Hauffä Miſt,\\
  in dem Sumär Wöderleuthen,\\
  döß erfordert gwiß än gſcheiten\\
  bin I nit ſelbſt mit darbey,\\
  gſchieht pur lautri Käierey.
\end{lyrics}

\clearpage
\begin{lyrics}{21 Recitativo}
  \voice[Mazko]
  Herr Richtär hamär ſcharffe Klagn\\
  wider Parthl vorzutragn,\\
  Schwabnwürth mus Vorſpröcher ſey,\\
  der woäß auß der Studi fei,\\
  wie mär Sachä her citirt\\
  und ſo änHandl recht außführt.

  \voice[Schwabnwürth]
  Wohl weiß und großgünſchtiger Herr Richter\\
  als ein kluger Fridens-Schlichter,\\
  ich wend jezſcht zwo Klägä ein\\
  vor dem Mazko und die Gmein,\\
  dem erſchte wolt er Orfeig gebä,\\
  auch thet er mit der gänzä Gmein änhöbä\\
  und hieſſ ſie den Hoblä ſauber außbläſä,\\
  diß riecht ein jedä in die Näſä.\\
  Nu wölle älle Männ vor Männ\\
  die grechte Satiſchfaction.

  \voice[Richter]
  Seyts zufrieden wann ihn leg\\
  in den Stokh auf zwoä drey Täg?

  \voice[Schwabnwürth]
  So khän ſchön die gäntze Gmein\\
  wie auch Mazko zufridä ſeyn.
\end{lyrics}

\begin{lyrics}{22 Chorus}
  \voice[alle]
  Mir habn än Richter oänen Mann\\
  der nach der Grechtikheit gethan,\\
  den wöll mär alle liebn und ehrn,\\
  ſo lang ſei Richter-Ambt thuet wern.\\
  Hiezt endli iſt d’Bohnthätung auß,\\
  än jeder geh zum Sterz nach Hauß.
\end{lyrics}

\end{document}
